% !TeX root = ../main.tex

\documentclass[../main.tex]{subfiles}
\begin{document}

\begin{center}
    \Large{\textbf{LỜI NÓI ĐẦU}}\\
\end{center}
\vspace{1cm}

Trong bối cảnh đô thị hóa và phát triển kinh tế diễn ra nhanh chóng, lượng chất thải rắn sinh hoạt phát sinh tại các đô thị Việt Nam ngày càng gia tăng về cả khối lượng lẫn mức độ phức tạp của thành phần. Nếu không được phân loại và xử lý phù hợp, chất thải rắn sinh hoạt không chỉ gây áp lực lên hệ thống thu gom, vận chuyển và xử lý mà còn làm phát sinh nhiều tác động bất lợi đến môi trường đất, nước, không khí và cảnh quan đô thị. Thực tiễn đó cho thấy nhu cầu cấp thiết phải chuyển từ mô hình xử lý thiên về chôn lấp sang các phương án có khả năng thu hồi tài nguyên và giảm phát thải.

Trong số các giải pháp hiện nay, công nghệ chuyển hóa chất thải thành năng lượng được xem là một hướng tiếp cận có ý nghĩa kỹ thuật và môi trường rõ rệt. Đối với những nguồn chất thải có thành phần cháy được tương đối lớn, công nghệ nhiệt phân kết hợp thu hồi nhiên liệu lỏng cho phép tận dụng phần hữu cơ có giá trị nhiệt, đồng thời giảm lượng chất thải phải chôn lấp và tạo ra dòng sản phẩm có khả năng lưu trữ, xử lý tiếp hoặc sử dụng cho phát điện. Đây là cơ sở để xem xét các phương án công nghệ phù hợp với điều kiện cụ thể của từng địa phương.

Xuất phát từ yêu cầu đó, đồ án tốt nghiệp này được thực hiện với đề tài \textit{Nghiên cứu công nghệ xử lý rác 120 tấn/ngày đêm, ứng dụng phát điện}. Nội dung đồ án tập trung vào việc phân tích đặc tính chất thải đầu vào, lựa chọn phương án công nghệ xử lý phù hợp, tính toán cân bằng vật chất và năng lượng ở mức thiết kế sơ bộ, đồng thời đề xuất các giải pháp vận hành, kiểm soát môi trường và an toàn công nghệ. Kết quả nghiên cứu hướng đến việc xây dựng một phương án kỹ thuật có tính khả thi, phù hợp với điều kiện thực tế của dự án và có thể làm cơ sở tham khảo cho các công trình tương tự.

Mặc dù đã cố gắng tiếp cận vấn đề một cách nghiêm túc, đồ án chắc chắn vẫn không tránh khỏi những thiếu sót do giới hạn về thời gian, kinh nghiệm và phạm vi dữ liệu. Em xin chân thành cảm ơn quý thầy cô đã tận tình hướng dẫn, đồng thời mong nhận được những ý kiến đóng góp quý báu để đồ án được hoàn thiện hơn.

\end{document}
